\section{Diario della riunione}
\subsection{Passaggio degli sprint da 2 settimane agli sprint da 1 settimana}
Il gruppo, con valutazione unanime, ha deciso di ridurre la durata degli sprint da 2 settimane ad 1 settimana, per facilitare la stesura di preventivi, consuntivi e una migliore organizzazione nei compiti assegnati.

\subsection{Cambio ruoli}
In questo sprint, i ruoli saranno:
\begin{itemize}
	\item Matteo Mazzaretto $\rightarrow$ Responsabile
	\item Luca Marcuzzo $\rightarrow$ Progettista
	\item Laura Pieripolli $\rightarrow$ Progettista
	\item Lisa Casagrande $\rightarrow$ Progettista
	\item Tommaso Ceron $\rightarrow$ Progettista
	\item Marco Brunello $\rightarrow$ Progettista
\end{itemize}


\subsection{Stesura consuntivo sprint precedente e preventivo sprint attuale}
Il gruppo ha segnato le ore produttive completate nello sprint precedente per redigere il consuntivo nel PdP.\\
Oltretutto il gruppo ha già steso le ore preventivate per questo sprint.

\subsection{Decisione punti di discussione con l'azienda}
I principali punti di discussione con l'azienda saranno:
\begin{itemize}
	\item Capire se è preferibile che l'applicazione avrà un sistema di deployment a monolite oppure a microservizi;
	\item Capire in quali punti sia ottimale usare i DB vettoriali;
	\item Inizio di progettazione architetturale, design thinking.
\end{itemize}

\subsection{Mostrare setup della repository di SmartOrder con varie indicazioni}
I due componenti del gruppo che avevano il compito di completare il setup iniziale della repository di SmartOrder hanno mostrato i loro progressi, evidenziando i punti di forza dei risultati raggiunti.

\subsection{Assegnazione compiti}
Si è deciso di assegnare il compito di modifica dei riferimenti normativi e informativi.\\
Oltretutto si è deciso di continuare per la strada di testare diversi modi di progettazione, attendendo l'incontro con l'azienda per definire in maniera più precisa i ruoli.

\subsection{Discussione stile di codifica}
Il gruppo ha deciso di avere uno stile di codifica unico, il quale sarà:
\begin{itemize}
	\item Commenti in italiano
	\item Nomenclatura di metodi, funzioni, classi e variabili in inglese
	\item Le classi saranno scritte in PascalCase
	\item Le variabili saranno scritte in camelCase
	\item I nomi delle funzioni saranno scritti in snake\_case
\end{itemize}